\scanpage{28}
\[
\operatorname{Coker}T_p(\widetilde\varphi)_{\text{ét}}\simeq(\operatorname{Ker}\unsure{(}\widetilde\varphi)(p)_{\text{ét}}
\]
Or on sait que $\operatorname{Ker}\widetilde\varphi$ est autodual, donc $(\operatorname{Ker}\widetilde\varphi)(p)_{\text{ét}}$ est dual de $(\operatorname{Ker}\widetilde\varphi(p))_{\mathrm{rad}}$, donc $=$ cardinal, d'où la relation
\[
\begin{array}{rcl}
\bigl(\operatorname{card}(\operatorname{Ker}\widetilde\varphi)(p)_{\text{ét}}\bigr)^2&=&\operatorname{card}\operatorname{Ker}(\widetilde\varphi(p))\\
&&\rotatebox{90}{$=$}\\
&& (\deg\widetilde\varphi)(p).
\end{array}
\]
d'où la relation envisagée \ldots\ldots\ . Cela définit bien un isom canonique
\diagram{$(*)$\quad\fbox{\begin{tikzcd}[ampersand replacement=\&,column sep=huge]
\det T_p(A)_{\text{ét}}\arrow[r,"\sim","\alpha_A"']\&\det T_p(A^*)_{\text{ét}}
\end{tikzcd}}}
compatible aux changt de base, et isom. en $A$.

[ Si $u$ est un autom. de $A$, et $\check u=u^{*-1}$ son contragredient, on aura
\[
\alpha\,\det T_p(u)_{\text{ét}}=\det T_p(\check u)_{\text{ét}}\alpha,
\]
donc en identifiant les dét s à des élém de $\mathbb Z_p$ \deleted{\unsure{(i)}}, on aura
\[
\begin{array}{rcl}
\det T_p(u)_{\text{ét}}&=&\det T_p(u^{*-1})_{\text{ét}}\\
&&\rotatebox{90}{$=$}\\
&& (\det T_p(u^*))^{-1}
\end{array}
\]
donc
\[
\boxed{\det T_p(u)_{\text{ét}}\,\inserted{\det}T_p(u^*)_{\text{ét}}=1}\qquad\text{\uline{$u$ un autom de $A$}}
\]
]

\margin{\figure{\begin{tikzpicture}[x=1cm,y=1cm]
\draw[line width=1pt] (0,1.4) .. controls (1.4,2) and (2.2,1.7) .. (1.65,1.2) .. controls (1.3,.9) and (.6,.55) .. (.4,.3) .. controls (.55,.1) and (1.35,.2) .. (1.9,.23);
\end{tikzpicture}}}
Attention, si $u$ est un endom de $A$, on aura en général
\[
\det T_p(u)\ne\det T_p(u^*)
\]
\mspage{27}
